\section{Characters and regularization on affine cones}\label{sec:D}

We translate the projective results of Sections~\ref{sec:B} and~\ref{sec:C}
into statements about homogeneous functions on affine cones.  The distinction
between a primitive prime order and a rescaled valuation is essential in this
translation.

\subsection{Angular fields and the discrepancy normalization}

Let $\breve R$ be a finitely generated normal domain over $k$, equipped with a
positive torus grading and a commuting diagonal action of a finite cyclic
group $G$.  Positivity means that there is a linear functional strictly
positive on every nonzero occurring weight and that the degree-zero ring is
$k$.  Set $R=\breve R^G$ and $X=\Spec R$, and assume that $X$ is
$\Q$-Gorenstein and klt.  We use the effective torus acting on $R$, with
character lattice $M$ of rank $r_T>0$.  Write
\[
 C=\operatorname{cone}\{m:R_m\ne0\}\subset M_{\R},\qquad
 \sigma=C^\vee\subset\operatorname{Hom}(M,\R),\qquad
 \mathcal K=\Frac(R)^{\mathbb T}.
\]
Thus $C$ is a full-dimensional pointed rational polyhedral cone.  Choose a
multiplicative splitting $m\mapsto z^m$ of the homogeneous rational
functions.  Such a splitting exists because $M$ is free abelian, and it gives
\begin{equation}\label{eq:D-angular-field}
 \Frac(R)=\mathcal K(z^M).
\end{equation}

Let $h\in\breve R$ be a nonzero homogeneous function of $G$-character $\chi$
and torus weight $\mathbf u\in M_{\Q}$.  If $r_G$ is the order of $G$, set
$H=r_G^{-1}\divisor(h^{r_G})$ on $X$.  Assume that there is a torus-invariant
rational top form $\theta$ on $X$ such that
\begin{equation}\label{eq:D-log-form}
 K_X(\theta)+H=0,\qquad (X,H)\text{ is lc}.
\end{equation}
Evaluation of a semi-invariant function always means evaluation of an
invariant power, divided by its exponent.  In particular, writing
$h^{r_G}=f_hz^{r_G\mathbf u}$ gives
\begin{equation}\label{eq:D-height-evaluation}
 (P,y)(h)=\gamma(P)+\langle y,\mathbf u\rangle,
 \qquad \gamma=\frac1{r_G}\divisor(f_h).
\end{equation}

For a primitive divisorial valuation $P$ of $\mathcal K$, define
\begin{equation}\label{eq:D-polyhedron}
 \Delta_P=\left\{y\in\operatorname{Hom}(M,\R):
 P(f)+\langle y,m\rangle\ge0\text{ whenever }fz^m\in R\right\}.
\end{equation}
The inequalities for finitely many homogeneous algebra generators suffice.
The polyhedron is nonempty: adding a sufficiently large multiple of a
strictly positive grading satisfies all these inequalities.  For rational
$y\in\Delta_P$, the Gauss rule
\[
 (P,y)\left(\sum_m f_mz^m\right)
   =\min_{f_m\ne0}\bigl(P(f_m)+\langle y,m\rangle\bigr)
\]
defines a valuation nonnegative on $R$.  After a positive rational
rescaling it is a primitive divisorial valuation.  For trivial restriction
to $\mathcal K$, the corresponding Gauss valuations are the nonzero rational
pure weights $y\in\sigma$, with value $\langle y,m\rangle$ on $fz^m$.

Choose a basis $z_1,\ldots,z_{r_T}$ of the weight monomials.  Torus
invariance gives a rational top form $\theta_0$ on $\mathcal K$ such that
\[
 \theta=\theta_0\wedge d\log z_1\wedge\cdots\wedge d\log z_{r_T}.
\]
Here $\theta_0$ is a nonzero constant when $\mathcal K=k$.  For a primitive
$P$ put $a(P)=\nu_{\log,P}(\theta_0)$, where log order is the order of a top
form plus one.  Extend $a$ homogeneously to positive multiples of prime
orders.

\begin{proposition}[Cone dictionary]\label{prop:D-cone-dictionary}
In the above setting the following statements hold.
\begin{enumerate}[label=\textup{(\roman*)}]
\item The discrepancy of $(X,H)$ at $(P,y)$ is $a(P)$, and its discrepancy
at a pure weight is zero.  Thus $a(P)$ is a nonnegative integer for every
primitive $P$, and $\mathbf u\in\operatorname{int}C$.
\item For rational $\mu\in\operatorname{int}C$, the ray quotient $V_\mu$ is
normal and projective, with function field $\mathcal K$.  The order function
\begin{equation}\label{eq:D-polarization}
 \mathbf L(\mu)(P)=\min_{y\in\Delta_P}\langle y,\mu\rangle
\end{equation}
is the b-pullback of an ample rational divisor on $V_\mu$.  Consequently
\begin{equation}\label{eq:D-height-divisor}
 \mathbf D(P)=\min_{y\in\Delta_P}(P,y)(h)
             =\gamma(P)+\mathbf L(\mathbf u)(P)
\end{equation}
is the b-pullback of an effective ample rational divisor $D$ on
$V_{\mathbf u}$.
\item At every actual prime $P$ of $V_\mu$, a minimizer $y$ of
\eqref{eq:D-polarization} can be chosen so that $e(P,y)$ is a primitive
prime order on $X$, for some positive integer $e$.  For this choice,
\begin{equation}\label{eq:D-actual-prime}
 a(P)+(P,y)(h)=\frac1e,
 \qquad e y\in\operatorname{Hom}(M,\Z).
\end{equation}
On each $V_\mu$, $a$ is the discrepancy function of an effective ordinary
lc Calabi--Yau pair, and $a+\mathbf D$ is that of an effective generalized
klt Calabi--Yau pair with nef part $\mathbf D$.  Each $V_\mu$ is of Fano
type.  At an actual prime outside the floor $\{a=0\}$, the same minimizer
satisfies
\begin{equation}\label{eq:D-off-floor}
 a(P)=1,\qquad e=1,\qquad (P,y)(h)=0.
\end{equation}
\end{enumerate}
\end{proposition}

\begin{proof}
Take a uniformizer $t$ at $P$ and separating residue coordinates
$s_1,\ldots,s_{p-1}$, where $p=\trdeg_k\mathcal K$.  Up to a unit the base
form is
$t^{a(P)}d\log t\wedge ds_1\wedge\cdots\wedge ds_{p-1}$.
For a rational Gauss extension, integral powers of the $z_j$ can be
adjusted by powers of $t$ to obtain angular units with algebraically
independent residues.  The resulting logarithmic wedge has log order zero.
Together with \eqref{eq:D-log-form}, this proves the discrepancy formula.
The calculation with trivial restriction to $\mathcal K$ is the same,
using a basis of the kernel of the pure weight.  Nonnegativity follows
from lc.  Integrality is the ordinary integral log order of $\theta_0$.
At a nonzero rational pure weight $y$, klt of $X$ gives
$A_X(y)=y(h)=\langle y,\mathbf u\rangle>0$.  Positivity on the nonzero
rays of $\sigma$ is equivalent to $\mathbf u\in\operatorname{int}C$.

Choose a positive integer $d$ with $d\mu\in M$ and set
\[
 V_\mu=\Proj\bigoplus_{j\ge0}R_{jd\mu}.
\]
The ray ring is finitely generated and normal: it is the appropriate
diagonalizable invariant subring, followed, if necessary, by a Veronese
subring.  Its degree-zero part is $k$, so its Proj is projective.  An
invariant rational function is a ratio of two homogeneous elements of the
same weight.  Multiplying their common weight by a suitable occurring
weight moves both elements onto a sufficiently divisible multiple of the
interior ray.  To see the latter assertion directly, express $\mu$ as a
strictly positive rational combination of all the finitely many generating
weights and clear denominators; sufficiently high multiples leave room to
subtract any prescribed occurring weight.  This proves
$k(V_\mu)=\mathcal K$.

Write the algebra generators as $f_jz^{m_j}$.  Linear programming duality
gives
\begin{equation}\label{eq:D-linear-program}
 \min_{y\in\Delta_P}\langle y,\mu\rangle
 =\max\left\{-\sum_j\lambda_jP(f_j):
       \lambda_j\ge0,\ \sum_j\lambda_jm_j=\mu\right\}.
\end{equation}
The polytope on the right is compact because the grading is positive.
Its rational vertices depend on the weights and $\mu$, not on $P$.
After clearing their denominators and taking a Veronese generated in
degree one, these vertices give actual monomials on the ray.  Conversely,
a sum of monomials has no worse pole order than its worst summand.
Therefore \eqref{eq:D-linear-program} is exactly the maximum allowed pole
order in the homogeneous linear system, expressed using $z^{d\mu}$ as
rational frame.  On the Proj this system is basepoint-free in sufficiently
divisible degree.  Its maximum-pole formula on every higher model is the
pullback of its ample tautological divisor.  This proves (ii), including
the b-divisor assertion.  Formula \eqref{eq:D-height-evaluation} identifies
$\mathbf D$ with a rational principal adjustment of this polarization.
Since $h$ is regular upstairs, its invariant powers have nonnegative
value at every Gauss extension, so $\mathbf D\ge0$.

For (iii), choose a homogeneous section $s$ on the defining ray whose
nonvanishing chart contains the generic point of $P$.  The map
$\Spec R_s\to\Spec(R_s)_0$ is the torus quotient, and the latter is the
corresponding chart of $V_\mu$.  Localize the base at $P$.  Its uniformizer
remains a nonunit upstairs: the invariant ring is a direct summand, and
an inverse upstairs would give an inverse downstairs by taking its
invariant part.  A component of its zero divisor therefore has height
one and dominates $P$, by the principal ideal theorem.  Such a component
is torus invariant, since a connected torus cannot permute its finitely
many irreducible components nontrivially.

The primitive order of a torus-invariant prime is Gauss.  Indeed its
valuation ideals are torus invariant, hence are direct sums of weight
spaces; there is no cancellation between different weights, and on each
weight space restriction to $\mathcal K$ supplies the indicated base
order.  Write the resulting valuation as $e(P,y)$.  Its restriction to
$\mathcal K$ is an integral multiple $eP$, and its values on $z^M$ are
integral.  Invertibility of $s$ implies that $y$ minimizes the ray
polarization.  A prime on a normal variety has ordinary log discrepancy
one.  Thus
$1=e\bigl(a(P)+(P,y)(h)\bigr)$, which proves
\eqref{eq:D-actual-prime}.  Notice that this argument concerns actual
primes on $V_\mu$; it asserts no such identity at an arbitrary exceptional
prime.

On $V_\mu$ set $\Omega=-K_{V_\mu}(\theta_0)$.  The log discrepancies of
$(V_\mu,\Omega)$ are $a$, and $K_{V_\mu}(\theta_0)+\Omega=0$.
At actual primes \eqref{eq:D-actual-prime} gives $0\le a\le1$; hence
$\Omega$ is effective.  Moreover
$a(P)+\mathbf D(P)\le1$ there.  At every divisorial $P$, a rational
minimizer for \eqref{eq:D-height-divisor} gives a Gauss extension, so klt
of $X$ implies $a(P)+\mathbf D(P)>0$.  Subtracting the trace of
$\mathbf D$ from $\Omega$ and taking $\mathbf D$ as nef part gives the
claimed generalized klt Calabi--Yau data.

For completeness these data imply Fano type without any uniform index
assumption.  On a common resolution on which $\mathbf D$ is semiample,
replace it by a general divided member of a sufficiently divisible free
system.  Added to the generalized boundary this gives an ordinary klt
sub-pair; its pushdown has effective boundary and is crepant, with total log
canonical divisor rationally trivial.  The added divisor is big.
Choose an effective ample divisor $A$ on $V_\mu$.  Bigness permits a second
effective representative of $\mathbf D$ containing a small positive
multiple of the pullback of $A$.  Mixing this representative with the
general free representative with sufficiently small positive weight
preserves klt.  The resulting effective boundary downstairs contains
$\varepsilon A$ for some $\varepsilon>0$.  Subtracting that summand leaves
an effective klt boundary whose negative log canonical class is
$\Q$-linearly equivalent to $\varepsilon A$, hence ample.  Finally, at
a non-floor prime the integer $a(P)>0$ in
\eqref{eq:D-actual-prime} forces \eqref{eq:D-off-floor}.
\end{proof}

\subsection{Which divisorial lc tests occur}

\begin{lemma}[All divisorial lc tests are Gauss]\label{lem:D-lc-tests}
Every nontrivial divisorial valuation nonnegative on $R$ and having
zero discrepancy for $(X,H)$ is Gauss.  Its restriction to $\mathcal K$
is either $eP$ with $a(P)=0$, or trivial, in which case it is a nonzero
pure weight in $\sigma$.  Conversely every rational Gauss extension
with $a(P)=0$, and every nonzero rational pure weight, has zero
discrepancy and a center whose closure contains the contracting vertex.

If $\Spec\breve R\to X$ is quasi-\'{e}tale and the pairs are related by
log pullback, every divisorial $h$-lc test upstairs restricts to one of
these tests, up to scaling.  Conversely all the rational tests just
listed lift to divisorial $h$-lc tests upstairs whose center closures
contain the upstairs vertex.  These statements suffice in particular
for testing invariant rational functions.
\end{lemma}

\begin{proof}
Let $v$ be a primitive divisorial valuation as in the first assertion.
By the divisorial restriction statement in Section~\ref{sec:A}, a
nontrivial restriction to $\mathcal K$ has the form $eP$, with $P$
primitive and $e$ a positive integer.  Use the uniformizer $t$ and
separating residue coordinates from the proof of
Proposition~\ref{prop:D-cone-dictionary}.  Every remaining logarithmic
differential has nonnegative log order at $v$.  Thus
\begin{equation}\label{eq:D-log-defect}
 A_{X,H}(v)\ge e a(P).
\end{equation}
More explicitly, pass to the finite-index sublattice of $M$ on which
$v(z^m)$ is divisible by $e$, and replace each of its basis monomials
by the unit $z^m t^{-v(z^m)/e}$.  If their residues are algebraically
dependent over the residue field of $P$, their leading differential
wedge vanishes, so the inequality in \eqref{eq:D-log-defect} is strict.
Zero discrepancy therefore forces $a(P)=0$ and independence of these
angular unit residues.  For a finite sum of tied terms, divide by a
term of least value.  All differences of its exponents lie in that
finite-index lattice; independence of the adjusted residues prevents
cancellation.  Thus $v$ is exactly the Gauss valuation determined by
its restrictions to $\mathcal K$ and $z^M$.

If the restriction to $\mathcal K$ is trivial, use a separating
transcendence basis of $\mathcal K$ in the same wedge computation.
If all angular weights were zero, the unit residues could have
transcendence degree at most $r_T-1$ over $\mathcal K$, and the leading
top wedge would vanish.  Zero log discrepancy therefore forces a
nonzero angular weight.  It further forces algebraic independence,
over $\mathcal K$, of the residues of a basis of the weight-zero
monomials.  The same tied-sum argument now proves that $v$ is a pure
Gauss weight.  Nonnegativity on $R$ places that weight in $\sigma$.

The converse discrepancy assertion was computed in
Proposition~\ref{prop:D-cone-dictionary}.  The center of a Gauss
valuation is torus invariant.  A positive grading contracts its
closure to the vertex, so the vertex lies in that closure; the
valuation need not itself be centered at the closed vertex.

For a finite quasi-\'{e}tale map with log-pulled pairs, the ramification
formula makes discrepancies homogeneous under restriction.  Hence
an upstairs lc valuation restricts to a downstairs lc valuation,
and every divisorial extension of a downstairs lc valuation is lc
upstairs.  Such extensions exist by taking the normalization of a
model carrying the downstairs prime.  The finite image of an upstairs
center closure is the downstairs center closure and contains the
vertex.  There is only one point upstairs over that vertex: every
positive homogeneous semi-invariant generator has an invariant
positive power, which vanishes there.  Therefore the upstairs center
closure contains the upstairs vertex as claimed.
\end{proof}

\subsection{Sections and a uniform sandwich argument}

Let $\xi=fz^k\in\Frac(R)$ be a nonzero rational eigenfunction, invariant
under the finite group, and let $m\in\Z$ satisfy
$\mu=k+m\mathbf u\in\operatorname{int}C$.  Define
\begin{equation}\label{eq:D-section-divisor}
 \mathbf J_m(P)=\min_{y\in\Delta_P}(P,y)(\xi h^m)
   =\divisor(f)(P)+m\gamma(P)+\mathbf L(k+m\mathbf u)(P).
\end{equation}
This is the b-pullback of an ample rational divisor $J_m$ on $V_\mu$;
the divisor itself may be signed.  At a non-floor prime of $V_\mu$ it has
integral coefficient, by \eqref{eq:D-off-floor} and the integral value of
the finite-group invariant function $\xi$ at a primitive cone prime.

For a rational b-divisor $\mathbf A$ that descends to a model, write
\[
 S(\mathbf A)=\{0\}\cup
 \{g\in\mathcal K^*: \divisor(g)+\mathbf A\ge0\}.
\]
Testing on its determining model is equivalent to testing all divisorial
orders.  This notation agrees with that of Section~\ref{sec:C} and also
applies to signed divisors.

\begin{lemma}[Sections as regular symbols]\label{lem:D-section-symbols}
For every positive integer $q$ and every $g\in\mathcal K$,
\begin{equation}\label{eq:D-symbol-equivalence}
 g\in S(q\mathbf J_m)
 \quad\Longleftrightarrow\quad
 g\xi^q h^{qm}\in\breve R.
\end{equation}
In particular the neutral functions of cost $b\in\N$ are precisely
\begin{equation}\label{eq:D-neutral-functions}
 \{g\in\mathcal K:g h^b\in\breve R\}=S(b\mathbf D).
\end{equation}
\end{lemma}

\begin{proof}
Regularity implies all Gauss inequalities, hence the implication to the
left.  Conversely take a sufficiently divisible power so that the
semi-invariant becomes invariant, its weight is on the integral ray, and
its degree is large enough for the ray ring to agree with the section ring
of the tautological polarization in that degree.  The divisor inequalities
and \eqref{eq:D-linear-program} then put that power in $R$.  The original
function belongs to $\Frac(\breve R)$ and is integral over $\breve R$, so
normality puts it in $\breve R$.  The second formula is the case $\xi=1$.
When $\mathcal K=k$, the same argument says that an interior weight becomes
effective in a divisible degree and then uses normality.
\end{proof}

All statements that follow are sequence statements in bounded dimension.
Constants implicit in $O(\cdot)$ are uniform in the sequence, and
hypotheses about small divisorial orders are imposed along every
subsequence, as in Sections~\ref{sec:A}--\ref{sec:C}.

\begin{proposition}[Sandwich regularization]\label{prop:D-sandwich}
Let $L\ge1$ vary with the sequence.  Choose integers $m>m_0$ with
$m-m_0\ge cL$ for a fixed $c>0$, and suppose both
$k+m_0\mathbf u$ and $k+m\mathbf u$ are interior.  Assume that whenever
primitive divisorial valuations $P$ satisfy $a(P)=0$ and
$L\mathbf D(P)\to0$, there are integers $j_P$ such that
\begin{equation}\label{eq:D-sandwich-hypothesis}
 j_P\le\mathbf J_{m_0}(P)\le\mathbf J_m(P)
       \le j_P+O(L\mathbf D(P)).
\end{equation}
Then there is a bounded positive integer $q$ and a nonzero
$g\in S(q\mathbf J_m)$.  Equivalently, $g\xi^qh^{qm}$ is regular on
$\Spec\breve R$.
\end{proposition}

\begin{proof}
For $\mathcal K=k$ the preceding lemma gives the assertion with $q=1$.
Otherwise work first on $V_\mu$, where $\mu=k+m\mathbf u$.  Use the
generalized klt Calabi--Yau discrepancy $a+\mathbf D$ to extract onto a
$\Q$-factorial Fano type model $U$ all exceptional primes below a fixed
small positive discrepancy cutoff, and no other exceptional primes.
The extraction statement is supplied by Section~\ref{sec:A}.
All extracted primes are floor primes: $a$ is integral and nonnegative.
On $U$, the positive coefficients of the trace $D_U$ are precisely the
floor primes, and they are at most one.  This holds for surviving primes
by \eqref{eq:D-actual-prime}, and for extracted primes by their small
$a+\mathbf D$ discrepancy.

For each floor prime $Q$ on $U$ put
$\rho_Q=\min\{1,L\mathbf D(Q)\}$.  Choose a small fixed rational
$c_1>0$, with $2c_1\le c$, and an effective correction $R_c$, supported
on these primes, such that
\begin{equation}\label{eq:D-rounding}
 c_1\rho_Q\le R_c(Q)\le2c_1\rho_Q,
 \qquad
 \operatorname{den}\bigl(J_m(Q)-R_c(Q)\bigr)\le C_1/\rho_Q.
\end{equation}
Here $\operatorname{den}$ is the reduced positive denominator.  Such a
choice is obtained by taking a rational grid of mesh at most
$c_1\rho_Q$ and denominator at most $C_1/\rho_Q$ in the interval
$[J_m(Q)-2c_1\rho_Q,J_m(Q)-c_1\rho_Q]$.  There is no correction outside
the floor: no such primes were extracted, and $J_m$ is integral there.

Minima of sums give the b-divisor inequality
\begin{equation}\label{eq:D-height-difference}
 \mathbf J_m-\mathbf J_{m_0}\ge(m-m_0)\mathbf D.
\end{equation}
Since $R_c\le2c_1 L D_U$, the divisor
$P_U=(J_m)_U-R_c$ dominates the trace of $\mathbf J_{m_0}$ and is big.
Run its MMP on $U$ to a nef big semiample divisor $P'$ on a
$\Q$-factorial Fano type model $U'$.  On a common higher model, denote
pullback from a displayed model by a star.  Then
\begin{equation}\label{eq:D-mmp-sandwich}
 \mathbf J_m\ge(P')^*\ge\mathbf J_{m_0},\qquad
 R_c^*\le \mathbf E:=\mathbf J_m-(P')^*\le(R_c')^*.
\end{equation}
Here is justification for both sides, including the comparison with the
possibly different ray $m_0$.  The divisor $J_m$ descends on $U$, and
MMP negativity gives $(P')^*\le P_U^*=\mathbf J_m-R_c^*$.
For the middle inequality, every divisible section of
$\mathbf J_{m_0}$ is a section of $P_U$, hence of $P'$ by the MMP.
The maximum-pole description of semiample divisors, applied on a common
determination, gives $(P')^*\ge\mathbf J_{m_0}$.  Finally b-nef
negativity gives $\mathbf J_m\le(J_m')^*$; since
$J_m'=P'+R_c'$, this is the last inequality in
\eqref{eq:D-mmp-sandwich}.  Also $\rho_Q\ge D_U(Q)$, so
\[
 \mathbf E\ge R_c^*\ge c_1(D_U)^*\ge c_1\mathbf D.
\]
The final comparison is again b-nef negativity for $\mathbf D$.

We apply Lemma~\ref{lem:B-phase-nonvanishing} to the polarization $P'$.
Choose fixed rational $\delta_2>0$ small enough for the denominator bound
below, and then $0<\delta_1\le\delta_2c_1$.  For
$0\le t\le\delta_2$ use the discrepancy family and nef part
\begin{align*}
 a_t&=a+(1-t/\delta_2)\delta_1\mathbf D+t\mathbf E,\\
 \mathbf M_t&=(1-t/\delta_2)\delta_1\mathbf D+t\mathbf J_m.
\end{align*}
Their total log canonical divisor is $\Q$-linearly equivalent to $tP'$.
For $t\ge\delta_2$ keep $a_t=a+\delta_2\mathbf E$ and add
$(t-\delta_2)(P')^*$ to the nef part.  This gives families on arbitrarily
long parameter intervals.  They are nondecreasing, klt, and have nef
$\Q$-Cartier b-parts.

We check effectivity and the denominator trace requirement of that lemma.
Every prime of $U'$ survives from $U$.  At a non-floor prime,
$a=1$, $\mathbf D=\mathbf E=0$, and $P'$ has integral coefficient.
At a surviving floor prime $Q$,
$\mathbf E(Q)=R_c(Q)\le2c_1\rho_Q$, while
\[
 \frac1{\operatorname{den}(P'(Q))}\ge\frac{\rho_Q}{C_1}.
\]
Take $2\delta_2c_1C_1\le1$.  Monotonicity gives
$0<a_t(Q)\le\delta_2R_c(Q)$ throughout the family, so
$a_t(Q)\le1/\operatorname{den}(P'(Q))$.  Equivalently, its effective
boundary dominates the denominator boundary $R(P')$ used in
Section~\ref{sec:B}.

It remains to check the signed phase hypothesis.  At a growing parameter,
small discrepancies mean $a(P)=0$ and $\mathbf E(P)\to0$.
Because $\mathbf E\ge c_1\mathbf D$, these valuations eventually belong
to the fixed-cutoff list already extracted on $U$.  On that list,
\eqref{eq:D-rounding} and \eqref{eq:D-mmp-sandwich} give
\[
 \mathbf E(P)\ge c_1\min\{1,L\mathbf D(P)\}.
\]
Thus $L\mathbf D(P)\to0$ and
$L\mathbf D(P)=O(\mathbf E(P))$.  Combining
\eqref{eq:D-sandwich-hypothesis} with
\eqref{eq:D-mmp-sandwich} yields
\[
 j_P\le(P')^*(P)\le j_P+O(\mathbf E(P)).
\]
The phase of $(P')^*(P)$ therefore has a nonnegative lift of size
$O(a_t(P))$, which is stronger than the signed phase hypothesis.
Lemma~\ref{lem:B-phase-nonvanishing} supplies a nonzero section in a
bounded positive degree of $P'$.  By
\eqref{eq:D-mmp-sandwich} it lies in $S(q\mathbf J_m)$, and
Lemma~\ref{lem:D-section-symbols} finishes the proof.
\end{proof}

\subsection{Bounded multiples, character counts, and tilted heights}

The sandwich proposition turns a phase interval into a regular symbol.
We use it in three forms: to supply symbols for bounded-width lattice
characters, to count symbols with one character and weight, and to
supply a symbol whose cost reflects a small tilted height.  The last form
will be used on newly optimized cones before any further rank refinement.

\begin{corollary}[Bounded-multiple character regularization]
\label{cor:D-regularization}
Suppose that the order condition \eqref{eq:C-alpha} on $V_{\mathbf u}$,
with error $\eta_i\to0$, holds whenever $a(P)=0$ and
$L\mathbf D(P)\to0$.  Fix $C_0>0$.  If $k\in M$ satisfies
\begin{equation}\label{eq:D-width}
 |\langle y,k\rangle|\le C_0L\langle y,\mathbf u\rangle
 \qquad(y\in\sigma),
\end{equation}
there are a bounded positive integer $q$, a nonzero invariant rational
eigenfunction $Y=fz^{qk}$, and a positive integer $b\le C_2L$ such that
$h^bY\in\breve R$.  The constants are uniform in $k$.  Both signs of
$k$ may be treated with the same bounded multiplier, and their costs may
be increased to a common cost $O(L)$.

The same bounds hold uniformly if $L$ is replaced by any $L'\ge L$ and
\eqref{eq:D-width} is imposed at scale $L'$.
\end{corollary}

\begin{proof}
Choose an integer $m_1\asymp L$ with $m_1>C_0L$, so that both
$\pm k+m_1\mathbf u$ lie in $\operatorname{int}C$.  Bigness of their
ray polarizations supplies rational-power functions $g_+,g_-$ on
$\mathcal K$ such that
\[
 (P,y)(z^{\pm k}h^{m_1})+P(g_\pm)\ge0
 \quad\text{for every }P\text{ and }y\in\Delta_P.
\]
Here a rational-power function means a positive integral root of an
element of $\mathcal K^*$ used only through its $\Q$-principal order;
it need not be an element of $\mathcal K$.  Put
\[
 \ell=-\divisor(g_+),\qquad
 G=2m_1D+\divisor(g_+g_-).
\]
Adding the two Gauss inequalities and taking the height minimum shows
that $G\ge0$ on $V_{\mathbf u}$, and
$G\sim_{\Q}2m_1D$.  At a non-floor prime $Q$ of this model, choose the
integral minimizing $y$ in \eqref{eq:D-off-floor}.  Since its height is
zero, the two inequalities give
\[
 0\le\langle y,k\rangle-\ell(Q)\le G(Q),
 \qquad \langle y,k\rangle\in\Z.
\]
Lemma~\ref{lem:C-scaled-phase} therefore gives a fixed positive integer
$q_0$ and a fixed $A>0$ such that, on all tiny sequences in the
statement, one can choose integers $j_P$ with
\[
 |q_0\ell(P)-j_P|\le A L\mathbf D(P).
\]

Apply Proposition~\ref{prop:D-sandwich} to $\xi=z^{q_0k}$.
For any sufficiently large integer $m_0=O(L)$ and $m=2m_0$, the first
Gauss inequality gives
\[
 \mathbf J_{m_0}\ge q_0\ell+(m_0-q_0m_1)\mathbf D.
\]
At a height-minimizing angular vector the second inequality gives
\[
 \mathbf J_m\le q_0\ell+(m-q_0m_1)\mathbf D+q_0\mathbf G
 \le q_0\ell+(m+q_0m_1)\mathbf D+q_0\mathbf G.
\]
The order condition \eqref{eq:C-alpha}, applied to
$G/(2m_1)\sim_{\Q}D$, bounds
$\mathbf G(P)\le2m_1(1+\eta_i)\mathbf D(P)$ on the tiny sequences.
Choose $m_0-q_0m_1\ge AL$.  These estimates give exactly
\eqref{eq:D-sandwich-hypothesis}.  A bounded further multiplier $q_1$
then supplies $g z^{q_1q_0k}h^{q_1m}\in\breve R$.  This is the desired
regularization.

Apply the argument to $-k$ as well, raise the two resulting functions
to bounded powers to align their multipliers, and multiply by powers of
$h$ to align their costs.  Uniformity in the weight follows by applying
the sequence argument to any proposed sequence of exceptions.  For an
enlarged scale $L'\ge L$, its tiny sequences satisfy
$L\mathbf D(P)\to0$ as well.  All the preceding hypotheses therefore
hold with $L'$ in place of $L$.  A sequence of exceptional pairs
$(L',k)$ would contradict the same proof; this proves the stated
uniformity, including for growing $L'/L$.
\end{proof}

\begin{corollary}[Count for a single character and weight]
\label{cor:D-character-count}
Assume the hypotheses of Lemma~\ref{lem:C-saturation-count} on
$V_{\mathbf u}$, and use its notation $p$ and $v_i$.  For a fixed
$C_0>0$, let $W_{t,k}$ be the space of regular symbols in $\breve R$
of finite-group character $\chi^t$ and torus weight $t\mathbf u+k$.
For every fixed positive integer $n$, eventually
\begin{equation}\label{eq:D-character-bound}
 \dim_k W_{t,k}\le C_3 v_i n^p
\end{equation}
simultaneously for all integers $0\le t\le C_0nL$ and all weights
satisfying
\[
 |\langle y,k\rangle|\le C_0nL\langle y,\mathbf u\rangle
 \qquad(y\in\sigma).
\]
The constant $C_3$ is independent of $n,t,k$.
\end{corollary}

\begin{proof}
If $W_{t,k}\ne0$, divide any nonzero member by $h^t$.  This is a
finite-group invariant rational eigenfunction, so $k\in M$.
Let $\mathbf J^{\mathrm{test}}$ denote the integral basepoint-free
test b-divisor called $J$ in Lemma~\ref{lem:C-saturation-count}.
That lemma has
$0\le\mathbf J^{\mathrm{test}}\le C L_0\mathbf D$ with
$L_0=O(L)$.  At a primitive $P$ with $a(P)=0$ and
$L\mathbf D(P)\to0$, its nonnegative integral value is eventually
zero.  Such $P$ belong to that lemma's test set $\mathcal L$ and hence
satisfy \eqref{eq:C-alpha}.  We may therefore apply
Corollary~\ref{cor:D-regularization} at the enlarged scale $nL$.

Choose an opposite invariant eigenfunction $Y$ of weight $-qk$, with
$q$ bounded and $h^bY$ regular for $b\le C_4 nL$.  For
$F\in W_{t,k}$ the neutral function
\[
 \Phi(F)=\frac{F^qY}{h^{qt}}
\]
satisfies $h^{qt+b}\Phi(F)=F^q h^bY\in\breve R$.  By
\eqref{eq:D-neutral-functions}, it belongs to
$S((qt+b)\mathbf D)\subseteq S(C_5nL\mathbf D)$.
Although $\Phi$ need not be linear, it is a homogeneous polynomial map
of degree $q$, nonzero away from the origin, and induces an injective
morphism
\[
 \mathbb P(W_{t,k})\longrightarrow
 \mathbb P\bigl(S(C_5nL\mathbf D)\bigr).
\]
Indeed proportional images imply that $(F_1/F_2)^q\in k^*$;
algebraic closedness of $k$ then implies $F_1/F_2\in k^*$.
An injective morphism of projective varieties preserves the dimension
of its source in its image, so
$\dim W_{t,k}\le\dim S(C_5nL\mathbf D)$.
The upper count in Lemma~\ref{lem:C-saturation-count}, used at an
integer at most $(C_5+1)n$, proves
\eqref{eq:D-character-bound}; $p$ is bounded by the ambient dimension.
The uniform enlarged-scale assertion in
Corollary~\ref{cor:D-regularization} makes $q,C_4,C_5$ independent of
$n,t,k$.  The eventual assertion retains the quantifier that $n$ is
fixed before passing sufficiently far along the sequence.
\end{proof}

\begin{corollary}[Characters for a tilted height]
\label{cor:D-tilted-characters}
Suppose there are boundedly many invariant rational eigenfunctions
$Y_j^\pm$ of weights $\pm k_j$, with
\[
 h^{m_j}Y_j^\pm\in\breve R,\qquad
 0<m_j\le A L_j,\qquad 1\le L_j\le L,
\]
where $A$ is fixed.  Fix bounded real numbers $u_j$ and $\delta>0$.
Assume that, at every nonzero pure weight and at all Gauss extensions
over primitive $P$ with $a(P)=0$, $L\mathbf D(P)\to0$, the evaluations
satisfy
\begin{equation}\label{eq:D-tilt-hypothesis}
 y_j=v(Y_j^+)=-v(Y_j^-),\qquad
 \sum_j\frac{u_jy_j}{L_j}\le\delta v(h).
\end{equation}
It suffices to check these inequalities on rational vectors.  For
fixed $M_0$ sufficiently large compared with $1/\delta$, choose integers
\[
 r_j=M_0u_jL/L_j+O(1),
\]
with uniformly bounded rounding errors.  Then there is a regular
nonzero symbol of finite-group character $\chi^{qm}$ and torus weight
\begin{equation}\label{eq:D-tilt-weight}
 q\left(m\mathbf u-\sum_jr_jk_j\right),
 \qquad 0<m\le C_6M_0\delta L,
\end{equation}
for a bounded positive integer $q$.  The constant $C_6$ depends only on
the fixed bounds for the number of functions, their costs, the $u_j$,
and rounding errors; it is independent of $\delta$ and $M_0$.
The bound for $q$ may depend on the fixed choices $\delta,M_0$.
No order condition \eqref{eq:C-alpha} is assumed here.
\end{corollary}

\begin{proof}
Set $\xi=\prod_j(Y_j^+)^{-r_j}$, which has weight
$-\sum_jr_jk_j$.  Regularity and opposite evaluations imply
$|y_j|\le m_jv(h)$ at every test under consideration.  Hence, with a
fixed constant $C_7$ controlling the rounding errors,
\begin{equation}\label{eq:D-tilt-lower}
 \sum_jr_jy_j
 \le(M_0\delta+C_7)Lv(h).
\end{equation}
Choose $M_0\delta\ge2C_7+2$ and an integer
$m_0\asymp M_0\delta L$ strictly larger than
$(M_0\delta+C_7)L$, and set $m=2m_0$.  The constants in these choices
are independent of $\delta,M_0$.  At pure weights
\eqref{eq:D-tilt-lower} implies strict positivity of the weights of
$\xi h^{m_0}$ and $\xi h^m$ on $\sigma\setminus\{0\}$; both rays
are therefore interior.  At the angular tests it gives
$\mathbf J_{m_0}(P)\ge0$.

At a height-minimizing angular vector, the elementary estimate
\[
 \left|\sum_jr_jy_j\right|
 \le\sum_j|r_j|m_j\mathbf D(P)
 \le C_8(M_0+1)L\mathbf D(P)
\]
also holds.  Consequently
$0\le\mathbf J_{m_0}(P)\le\mathbf J_m(P)
\le C_9(M_0,\delta)L\mathbf D(P)$ on the tiny sequences.
Proposition~\ref{prop:D-sandwich} applies with $j_P=0$ and produces
$g\xi^qh^{qm}\in\breve R$ for bounded $q$.  Its character, weight,
and cost are exactly those in \eqref{eq:D-tilt-weight}.
\end{proof}

